\section{结论}

本文建立了一个可复现的纳米硒细胞毒性建模流程，包含数据审计、按文献分组的嵌套交叉验证、无泄漏预处理和结果产物归档。严格重建使用 434 条记录和 24 个文献组；无 IC50 的随机森林基础模型 OOF MAE 为 19.11\%、RMSE 为 24.12\%、$R^2=0.109$，XGBoost 分别为 19.49\%、23.97\% 和 0.120。

加入元素 Se 含量、稳定性、结合能 Se3d 和晶型指标没有带来稳定的误差降低；显式 log 剂量--log 时间交互项的点估计也变差，按文献组 bootstrap 区间跨越 0。这些阴性结果说明，当前归档数据不足以支持化学增强或交互项改善文献外泛化的结论。

General 到 SeNP 的迁移对照也没有显示稳定收益。XGBoost 迁移相对 SeNP-only 的 MAE 点估计下降 0.67 个百分点，但按文献组 bootstrap 区间跨越 0；随机森林迁移的 MAE 上升 0.58 个百分点。通用模型直接 zero-shot 预测在两个模型上都更差，说明在当前特征契约下，通用数据不能直接替代 SeNP 专用预测模型。

IC50 辅助 XGBoost 的严格指标在数值上最好（MAE 18.35\%、RMSE 23.25\%、$R^2=0.172$），但 IC50 与响应终点相关，且配对改进区间跨越 0。因此它只能作为敏感性参考，不能作为前瞻部署模型。外部验证仍需在湿实验前冻结预测、补齐材料表征，并收集不把实验 IC50 作为输入的独立剂量--反应数据；当前仓库还缺少修饰分类、分子量、配方标识和前瞻实验字段。
